\chapter[Considerações Finais]{CONSIDERAÇÕES FINAIS}

Ao olhar para minha trajetória completa na AGES, percebo que minha evolução profissional aconteceu de forma muito mais concreta do que eu imaginava quando entrei no curso. A Engenharia de Software me deu uma base importante, principalmente em requisitos, processos, organização de tarefas e entendimento do ciclo de desenvolvimento. Porém, foi na AGES que esses conceitos deixaram de ser apenas conteúdo de disciplina e passaram a aparecer em problemas reais, com clientes reais, limitações reais e pessoas trabalhando com diferentes níveis de experiência.

No início, eu via o desenvolvimento de software de uma forma muito mais limitada, quase como escrever código e entregar funcionalidades. Com o tempo, essa visão mudou bastante. Aprendi que um projeto envolve requisitos, comunicação, divisão de responsabilidades, decisões técnicas, validação com cliente e capacidade de lidar com imprevistos. Essa mudança de visão foi uma das maiores contribuições da AGES para minha formação.

Em termos de \textit{hard skills}, a AGES foi essencial porque foi onde dei meus primeiros passos como desenvolvedor \textit{frontend}, área na qual atuo atualmente. Tive contato prático com React Native, React, TypeScript, Next.js, integração com APIs, autenticação, banco de dados, Git e deploy. Mais do que aprender ferramentas isoladas, comecei a entender como essas partes se conectam dentro de um produto real.

Esse crescimento técnico também teve impacto direto na minha vida profissional. A experiência prática da AGES contribuiu para eu conseguir meu primeiro estágio na área, porque me deu vivência real para conversar sobre desenvolvimento, dificuldades, trabalho em equipe e entregas.

Nas \textit{soft skills}, minha evolução também foi grande. Aprendi a me comunicar melhor e, principalmente, a me comunicar cedo. Entendi que esconder bloqueios ou esperar demais para pedir ajuda prejudica o time inteiro. Também desenvolvi mais clareza para explicar problemas, orientar colegas e passar conhecimento adiante, especialmente quando precisei apoiar alunos com menos experiência.

Uma das principais lições que levo é que projetos de software dão certo quando existe alinhamento entre pessoas, processo e técnica. Quando as tarefas são claras, a comunicação é transparente e cada pessoa entende sua responsabilidade, o time consegue evoluir mesmo com dificuldades. Quando isso não acontece, até tarefas simples podem virar grandes problemas.

Minha visão sobre Engenharia de Software, portanto, mudou bastante. Hoje entendo que código continua sendo importante, mas é apenas uma parte do processo. Engenharia de Software envolve entender contexto, organizar trabalho, lidar com pessoas, tomar decisões e manter o produto alinhado com o que o cliente realmente precisa.

Particularmente, considero a AGES uma experiência extremamente valiosa dentro do curso. Ela é um espaço muito bom para os alunos darem o primeiro passo em projetos reais, errarem, aprenderem, lidar com clientes, colaborar com colegas e entender melhor como o desenvolvimento de software acontece fora de exercícios isolados. No meu caso, a AGES foi decisiva não só para aprender tecnologias, mas para amadurecer como profissional, descobrir melhor meu caminho no \textit{frontend} e entender que ser engenheiro de software exige tanto conhecimento técnico quanto responsabilidade, comunicação e capacidade de aprender continuamente.
